\par\Needspace{6\baselineskip}\medskip
\SourceLeaf{213}\hypertarget{native-213-article157}{}\pdfbookmark[2]{157.}{article-157}
\begin{center}157.\end{center}
We shall now first determine the positions of the great circles drawn from the heliocentric positions of the Earth to the geocentric positions of the planet. We suppose the characters $\mathfrak A,\mathfrak A',\mathfrak A''$ assigned to the intersections of these circles with the equator, or, if one prefers, to their ascending nodes, and denote the distances of the points $B,B',B''$ from those intersections by $\Delta,\Delta',\Delta''$. In most of the operations we must now replace $A,A',A''$ by $\mathfrak A,\mathfrak A',\mathfrak A''$, and $\delta,\delta',\delta''$ by $\Delta,\Delta',\Delta''$. The attentive reader will readily see, even without our pointing it out, where $A,A',A'',\delta,\delta',\delta''$ must be retained.

\Needspace{15\baselineskip}
Carrying out the computation, we now find
\begin{center}
\begingroup\footnotesize\setlength{\tabcolsep}{3pt}\renewcommand{\arraystretch}{1.25}
\begin{tabular}{l|r|r|r}
\shortstack{Right asc. of points\\$\mathfrak A,\mathfrak A',\mathfrak A''$} & $233^0\,54'\,57''{,}10$ & $253^0\,8'\,57''{,}01$ & $276^0\,40'\,25''{,}87$\\
$\gamma,\gamma',\gamma''$ & $51\quad17\quad15{,}74$ & $90\quad1\quad3{,}19$ & $131\quad59\quad58{,}03$\\
$\Delta,\Delta',\Delta''$ & $215\quad58\quad49{,}27$ & $212\quad52\quad48{,}96$ & $220\quad9\quad12{,}96$\\
$\delta,\delta',\delta''$ & $56\quad26\quad34{,}19$ & $55\quad26\quad31{,}79$ & $69\quad10\quad57{,}84$\\
$\mathfrak A'D,\mathfrak AD',\mathfrak AD''$ & $23\quad54\quad52{,}13$ & $30\quad18\quad3{,}25$ & $29\quad8\quad43{,}32$\\
$\mathfrak A''D,\mathfrak A''D',\mathfrak A'D''$ & $33\quad3\quad26{,}35$ & $31\quad59\quad21{,}14$ & $22\quad20\quad6{,}91$\\
$\epsilon,\epsilon',\epsilon''$ & $47\quad1\quad54{,}69$ & $89\quad34\quad57{,}07$ & $42\quad33\quad41{,}17$\\
logarithms of sines & $9{,}864\,3525$ & $9{,}999\,9885$ & $9{,}830\,1910$\\
$\log\sin\tfrac12\epsilon'$ & & $9{,}847\,8971$ &\\
$\log\cos\tfrac12\epsilon'$ & & $9{,}851\,0614$ &
\end{tabular}\endgroup
\end{center}

In the computation of Article 138, for $l'$ we shall use the right ascension of the point $\mathfrak A'$. Thus we find
\begin{center}\begin{tabular}{ll}
$\log T\sin t$ .\,.\,.\,. & $8{,}486\,8236_n$\\
$\log T\cos t$ .\,.\,. & $9{,}284\,8162_n$.
\end{tabular}\end{center}
\noindent Hence $t=189^0\,2'\,48''{,}83$, $\log T=9{,}290\,2527$. Further, we have $t+\gamma'=279^0\,3'\,52''{,}02$,
\begin{center}\begin{tabular}{ll}
$\log S$ .\,.\,.\,.\,.\,.\,.\,.\,.\,.\,. & $9{,}011\,0566_n$\\
$\log T\sin(t+\gamma')$ .\,.\,. & $9{,}284\,7950_n$,
\end{tabular}\end{center}
\noindent whence $\Delta'-\sigma=208^0\,1'\,55''{,}64$, and $\sigma=4^0\,50'\,53''{,}32$.

In the formulas of Article 140 for $a$, $b$ and $\dfrac b a$ we must retain $\sin\delta$, $\sin\delta'$, $\sin\delta''$ themselves, and likewise in the formulas of Article 142. For these computations we have
\SourceLeaf{214}
\begin{center}\begingroup\footnotesize\setlength{\tabcolsep}{4pt}\renewcommand{\arraystretch}{1.3}
\begin{tabular}{ll}
$\mathfrak A''D'-\Delta''=171^0\,50'\,8''{,}18;$ & \shortstack{$\log\sin$ .\,.\,. $9{,}152\,3306;$\\$\log\cos$ .\,. $9{,}995\,5759_n$}\\
$\mathfrak AD'-\Delta=174\quad19\quad13{,}98;$ & \shortstack{--- .\,.\,. $8{,}995\,4722;$\\--- .\,.\,. $9{,}997\,8629_n$}\\
$\mathfrak A''D-\Delta''=172\quad54\quad13{,}39;$ & --- .\,.\,. $9{,}091\,7972$\\
$\mathfrak A'D-\Delta'+\sigma=175\quad52\quad56{,}49;$ & --- .\,.\,. $8{,}856\,1520$\\
$\mathfrak AD''-\Delta=173\quad9\quad54{,}05;$ & --- .\,.\,. $9{,}075\,5844$\\
$\mathfrak A'D''-\Delta'+\sigma=174\quad18\quad11{,}27;$ & --- .\,.\,. $8{,}996\,7978$.
\end{tabular}\endgroup\end{center}
\noindent From this we obtain
\[
\begin{aligned}
\log\varkappa&=0{,}921\,1850&\log\lambda&=0{,}081\,2057_n\\
\log\varkappa''&=0{,}811\,2762&\log\lambda''&=0{,}031\,9691_n\\
\log a&=0{,}109\,9088,&a&=+1{,}287\,9790\\
\log b&=0{,}181\,0404&&\\
\log\frac b a&=0{,}071\,1314&&
\end{aligned}
\]
\noindent which gives $\log b=0{,}181\,0402$. Between these two nearly equal values we shall adopt the mean $\log b=0{,}181\,0403$. Finally we obtain
\[
\begin{aligned}
\log c&=1{,}045\,0295\\
d&=+0{,}448\,9906\\
\log e&=9{,}210\,2894,
\end{aligned}
\]
\noindent which completes the preliminary computations.

The time interval between the second and third observations is $39{,}874\,409$ days, and between the first and second, $30{,}900\,961$ days. This gives $\log\theta=9{,}836\,2757$, $\log\theta''=9{,}725\,5533$. We therefore set \textit{for the first hypothesis}
\[
\begin{aligned}
x=\log P&=9{,}889\,2776\\
y=\log Q&=9{,}561\,8290.
\end{aligned}
\]
\noindent The principal results of the computation are then as follows:
\[
\begin{aligned}
\omega+\sigma&=20^0\,8'\,46''{,}72\\
\log Qc\sin\omega&=0{,}028\,2028.
\end{aligned}
\]
\noindent This gives the true value $z=21^0\,11'\,24''{,}30$, and $\log r'=0{,}350\,9379$. In this case the three remaining values of $z$ satisfying equation IV of Article 141 are
\[
\begin{aligned}
z&=63^0\,40'\,50''\\
z&=101\quad12\quad58\\
z&=199\quad24\quad7,
\end{aligned}
\]
\noindent\SourceLeaf{215}of which the first is to be regarded as an approximation to the orbit of the Earth. Its deviation is of course much greater here than in the preceding example, because the interval of time is too long. --- Further computation gives the following numbers:
\begin{center}\begingroup\footnotesize\setlength{\tabcolsep}{6pt}\renewcommand{\arraystretch}{1.3}
\begin{tabular}{ll}
$\zeta=195^0\,12'\,2''{,}48$ & $\zeta''=196^0\,57'\,50''{,}78$\\
$\log r=0{,}364\,7022$ & $\log r''=0{,}335\,5758$\\
$\tfrac12(u''+u)=266^0\,47'\,50''{,}47$ & $\tfrac12(u''-u)=-43^0\,39'\,5''{,}33$
\end{tabular}\endgroup\end{center}
\[
\begin{aligned}
2f'&=22^0\,32'\,40''{,}86\\
2f&=13\quad5\quad41{,}17\\
2f''&=9\quad27\quad0{,}05.
\end{aligned}
\]
\noindent We shall apportion the difference between $2f'$ and $2f+2f''$, which here is $0''{,}36$, to $2f$ and $2f''$ by setting $2f=13^0\,5'\,40''{,}96$, $2f''=9^0\,26'\,59''{,}90$.

\Needspace{17\baselineskip}
The times must now be corrected for aberration. In the formulas of Article 145 we must put $AD'-\zeta=\mathfrak AD'-\Delta+\delta-\zeta$, $A''D'-\zeta''=\mathfrak A''D'-\Delta''+\delta''-\zeta''$. We therefore have
\begin{center}\begin{tabular}{ll}
$\log r$ .\,.\,.\,.\,.\,.\,.\,.\,.\,.\,.\,. & $0{,}36470$\\
$\log\sin(AD'-\zeta)$ .\,.\,. & $9{,}76462$\\
$\mathrm C.\log\sin\delta$ .\,.\,.\,.\,.\,.\,. & $0{,}07918$\\
$\log\mathrm{const.}$ .\,.\,.\,.\,.\,.\,.\,. & $7{,}75633$\\\hline
& $7{,}96483$\\
reduction of time\ $=$ & $0{,}009222$
\end{tabular}\end{center}
\begin{center}\begin{tabular}{ll}
$\log r'$ .\,.\,.\,.\,.\,. & $0{,}35094$\\
$\log\sin(\delta'-z)$ & $9{,}75038$\\
$\mathrm C.\log\sin\delta'$ . & $0{,}08431$\\
& $7{,}75633$\\\hline
& $7{,}94196$\\
& $0{,}008749$
\end{tabular}\end{center}
\begin{center}\begin{tabular}{ll}
$\log r''$ .\,.\,.\,.\,.\,.\,.\,.\,.\,. & $0{,}33558$\\
$\log\sin(A''D'-\zeta'')$ & $9{,}84220$\\
$\mathrm C.\log\sin\delta''$ .\,.\,.\,.\,. & $0{,}02932$\\
& $7{,}75633$\\\hline
& $7{,}96343$\\
& $0{,}009192$
\end{tabular}\end{center}
\noindent This gives
\begin{center}\begin{tabular}{r|r|r}
Corrected times & intervals & logarithms\\
Nov.\ $5{,}564\,852$ & &\\
& $30{,}901\,434$ & $1{,}489\,9785$\\
$36{,}466\,286$ & &\\
& $39{,}873\,966$ & $1{,}60\overset{.}{0}\,6894$\\
$76{,}340\,252$ & &
\end{tabular}\end{center}
\noindent from which we derive the corrected logarithms of the quantities $\theta,\theta''$, respectively $9{,}836\,2708$ and $9{,}725\,5599$. Beginning the computation of the elements from $r',r'',2f,\theta$ gives $\log\eta=0{,}003\,1921$, just as from $r,r',2f'',\theta''$ we obtain $\log\eta''=0{,}001\,7300$. This yields $\log P'=9{,}890\,7512$, $\log Q'=9{,}571\,2864$, and hence
\[X=+0{,}001\,4736,\qquad Y=+0{,}009\,4574.\]

The principal results of the \textit{second hypothesis}, in which we set
\SourceLeaf{216}
\[
\begin{aligned}
x=\log P&=9{,}890\,7512\\
y=\log Q&=9{,}571\,2864,
\end{aligned}
\]
\Needspace{14\baselineskip}
\noindent are as follows:
\begin{center}\begingroup\footnotesize\setlength{\tabcolsep}{6pt}\renewcommand{\arraystretch}{1.3}
\begin{tabular}{r|r}
$\omega+\sigma=20^0\,8'\,0''{,}87$ & $\log r''=0{,}336\,9708$\\
$\log Qc\sin\omega=0{,}037\,3071$ & $\tfrac12(u''+u)=267^0\,6'\,10''{,}75$\\
$z=21^0\,12'\,6''{,}09$ & $\tfrac12(u''-u)=-43\quad39\quad4{,}00$\\
$\log r'=0{,}350\,7110$ & $2f'=22\quad32\quad8{,}69$\\
$\zeta=195^0\,16'\,59''{,}90$ & $2f=13\quad1\quad54{,}65$\\
$\zeta''=196\quad52\quad40{,}63$ & $2f''=9\quad30\quad14{,}38$\\
$\log r=0{,}363\,0642$ &
\end{tabular}\endgroup\end{center}
\noindent The difference $0''{,}34$ between $2f'$ and $2f+2f''$ is to be divided so that $2f=13^0\,1'\,54''{,}45$, $2f''=9^0\,30'\,14''{,}24$.

If it seems worth recomputing the corrections of the times here, we shall find $0{,}009\,169$ for the first observation, $0{,}008\,742$ for the second, and $0{,}009\,236$ for the third, and hence the corrected times Nov.\ $5{,}564\,905$, Nov.\ $36{,}466\,293$, Nov.\ $76{,}340\,208$. This gives
\[
\begin{aligned}
\log\theta&=9{,}836\,2703,&\log\theta''&=9{,}725\,5594\\
\log\eta&=0{,}003\,1790,&\log\eta''&=0{,}001\,7413\\
\log P'&=9{,}890\,7268,&\log Q'&=9{,}571\,0593.
\end{aligned}
\]
\noindent In this way the second hypothesis therefore yields
\[X=-0{,}000\,0244,\qquad Y=-0{,}000\,2271.\]

Finally, in the \textit{third hypothesis}, in which we set
\[
\begin{aligned}
x=\log P&=9{,}890\,7268\\
y=\log Q&=9{,}571\,0593,
\end{aligned}
\]
\noindent the principal results of the computation are as follows:
\begin{center}\begingroup\footnotesize\setlength{\tabcolsep}{6pt}\renewcommand{\arraystretch}{1.3}
\begin{tabular}{r|r}
$\omega+\sigma=20^0\,8'\,1''{,}62$ & $\log r''=0{,}336\,9536$\\
$\log Qc\sin\omega=0{,}037\,0857$ & $\tfrac12(u''+u)=267^0\,5'\,53''{,}09$\\
$z=21^0\,12'\,4''{,}60$ & $\tfrac12(u''-u)=-43\quad39\quad4{,}19$\\
$\log r'=0{,}350\,7191$ & $2f'=22\quad32\quad7{,}67$\\
$\zeta=195^0\,16'\,54''{,}08$ & $2f=13\quad1\quad57{,}42$\\
$\zeta''=196\quad52\quad44{,}45$ & $2f''=9\quad30\quad10{,}55$\\
$\log r=0{,}363\,0960$ &
\end{tabular}\endgroup\end{center}
\noindent\SourceLeaf{217}The difference $0''{,}30$ will here be divided so that $2f=13^0\,1'\,57''{,}25$, $2f''=9^0\,30'\,10''{,}42$\footnote{This rather large difference, nearly equal in all the hypotheses, arose chiefly because $\sigma$ had been found \underline{n}early two hundredths of a second of arc too small and the logarithm of $b$ a few units too large.}.

Since all these numbers differ very little from those furnished by the second hypothesis, we may now safely conclude that the third hypothesis will require no further correction and that a new hypothesis is therefore superfluous. We may accordingly proceed now to compute the elements from $2f',\theta',r,r''$. Since the computation consists of operations already explained in full above, it will suffice to give the resulting elements here for the benefit of those who wish to carry it out on their own:
\begin{center}\begingroup\footnotesize\setlength{\tabcolsep}{3pt}\renewcommand{\arraystretch}{1.3}
\begin{tabular}{ll}
\shortstack[l]{Right ascension of the ascending\\node on the equator\quad .\,.\,.\,.\,.\,.\,.\,.\,.\,.\,.\,.\,.\,.\,.\,.} & $158^0\,40'\,38''{,}93$\\
Inclination of the orbit to the equator\quad .\,.\,.\,.\,.\,.\,.\,.\,.\,.\,.\,.\,.\,.\,.\,.\,.\,.\,.\,.\,.\,.\,.\,.\,.\,.\,. & $11\quad42\quad49{,}13$\\
\shortstack[l]{Distance of the perihelion\\from that ascending node\quad .\,.\,.\,.\,.\,.\,.\,.\,.\,.\,.\,.\,.\,.\,.\,.\,.\,.\,.} & $323\quad14\quad56{,}92$\\
Mean anomaly for epoch 1806\quad .\,.\,.\,.\,.\,.\,.\,.\,.\,.\,.\,.\,.\,.\,.\,.\,.\,.\,.\,.\,.\,.\,.\,.\,. & $335\quad4\quad13{,}05$\\
Mean (sidereal) daily motion\quad .\,.\,.\,.\,.\,.\,.\,.\,.\,.\,.\,.\,.\,.\,.\,.\,.\,.\,.\,.\,.\,.\,.\,.\,.\,.\,.\,.\,.\,.\,. & $770''{,}2663$\\
Angle $\varphi$\quad .\,.\,.\,.\,.\,.\,.\,.\,.\,.\,.\,.\,.\,.\,.\,.\,.\,.\,.\,.\,.\,.\,.\,.\,.\,.\,.\,.\,.\,.\,.\,.\,.\,.\,.\,.\,.\,.\,.\,.\,.\,.\,.\,.\,.\,. & $14^0\,9'\,3''{,}91$\\
Logarithm of the semimajor axis\quad .\,.\,.\,.\,.\,.\,.\,.\,.\,.\,.\,.\,.\,.\,.\,.\,.\,.\,.\,.\,.\,.\,.\,.\,.\,.\,.\,.\,.\,.\,. & $0{,}442\,2438$.
\end{tabular}\endgroup\end{center}
